\subsection{Data Definition Language (DDL)}
The DDL is used to describe the schema of relations, in \sql\ called tables.
For that, we provide a table name, the columns and their data types.

\sql\ supports the following data types (plus more):
\begin{multicols}{3}
    \begin{itemize}[itemsep=-2pt]
        \item \texttt{CHAR(n)}
        \item \texttt{VARCHAR(n)}
        \item \texttt{INT}
        \item \texttt{BIGINT}
        \item \texttt{SMALLINT}
        \item \texttt{FLOAT}
        \item \texttt{BLOB(n)} (for large binaries)
        \item \texttt{DATETIME}
        \item \texttt{MONEY}
    \end{itemize}
\end{multicols}


\subsubsection{Creating Tables}
Use the DDL \texttt{CREATE TABLE} keyword to create a table. We can set default values for certain attributes,
or can add constraints to them.

Since the primary key must be unique for each entry, it may be useful to configure the primary key attribute with the following constraints
\texttt{PrimaryKeyField INT NOT NULL AUTO\_INCREMENT}

\inputcodewithfilename{sql}{}{code/sql/ddl/create.sql}

Note that PostgreSQL doesn't support \texttt{AUTO\_INCREMENT} constraints, instead use the \texttt{SERIAL} (or \texttt{BIGSERIAL}) type.


% TODO: Make sure all the sql statements actually execute in pgsql
\subsubsection{Updating / Altering Tables}
Use the DDL \texttt{ALTER TABLE} keyword to update a table's schema.
Be aware that if we ADD a column to the schema, unless a default value is provided, the column is set to \texttt{NULL} (see \ref{sec:sql-null})

\inputcodewithfilename{sql}{}{code/sql/ddl/alter.sql}

Instead of a \texttt{ADD COLUMN} or \texttt{DROP COLUMN}, we can also use \texttt{ADD CONSTRAINT} and \texttt{DROP CONSTRAINT} to update constraints


\subsubsection{Deleting Tables}

\inputcodewithfilename{sql}{}{code/sql/ddl/delete.sql}


\subsubsection{Constraints}
\sql\ supports a number of constraints that can be put on the different attributes of a schema.
\begin{itemize}
    \item \texttt{NOT NULL}: Disallows this column to be set to \texttt{NULL}, or have a \texttt{NULL} value
    \item \texttt{UNIQUE}: The value of each value in this column must be unique, but it can be \texttt{NULL}.
          An arbitrary number of columns can have \texttt{NULL} and still be valid
    \item \texttt{PRIMARY KEY}: Sets this column as the primary key. \texttt{NOT NULL} and \texttt{UNIQUE} is set on it implicitly.
          Contrary to common intuition, it can be set on multiple columns.
    \item \texttt{CHECK c}: Check that the values fulfil a condition.
          This condition \texttt{c} has the same syntax as for the \texttt{WHERE} clause (see \ref{sec:sql-basic-ops})
    \item \texttt{REFERENCES table}: A Foreign Key, used to refer to another tuple from a different relation.
          It is typically referencing the primary key of the other table.
\end{itemize}
For \texttt{REFERENCES}, there are many options for maintenance, which can be set on creating a reference, using \texttt{ON UPDATE} and \texttt{ON DELETE}:
\begin{itemize}
    \item \bi{Cascade}: Propagate \texttt{UPDATE} and \texttt{DELETE}
    \item \bi{Restrict}: Prevent deletion of the primary key before the change, causes error
    \item \bi{No Action}: Prevent modifications after attempting change, causes error
    \item \bi{Set default / Set null}: Set references to \texttt{NULL} or default value
\end{itemize}
The default for \texttt{ON DELETE} is \texttt{NO ACTION}.

\inlineexample
\mint{sql}|CREATE TABLE tab (id integer REFERENCES OtherTable ON DELETE cascade ON UPDATE cascade)|

\inlineintuition For a deletion of a row with \texttt{NO ACTION} set (the default), if there exist any rows of tables with a Foreign key on this row,
where the Foreign Key is equal to the to be delted primary key, an error is thrown.
If \texttt{RESTRICT} is set, the error happens a earlier, as the check is performed \textit{before} the deletion is attempted.
This has the downside of being possibly a bit slower, if large numbers of unreferenced rows are deleted.
